\FigureLayoutDeclare{img/1964/national/q6_fig1.png}{8.0cm}{0bp 0bp 0bp 0bp}

\chapter{1964年全国卷}

\section{解答题}

\begin{problem}
化简：\(\frac{\sqrt[3]{3^{-\frac32}\left( \frac13 \right)^{-3}}}{\left( \sqrt3-1 \right)^2}\).
\end{problem}

\begin{solution}
先化简根号内的幂：
\[
3^{-\frac32}\left( \frac13 \right)^{-3}=3^{-\frac32}3^3=3^{\frac32}
\]
因此
\[
\frac{\sqrt[3]{3^{-\frac32}\left( \frac13 \right)^{-3}}}{\left( \sqrt3-1 \right)^2}
=\frac{3^{\frac12}}{\left( \sqrt3-1 \right)^2}
=\frac{\sqrt3}{4-2\sqrt3}
=\frac{\sqrt3}{2\left( 2-\sqrt3 \right)}
\]
分母有理化，得
\[
\frac{\sqrt3}{2\left( 2-\sqrt3 \right)}
=\frac{\sqrt3\left( 2+\sqrt3 \right)}{2\left( 2-\sqrt3 \right)\left( 2+\sqrt3 \right)}
=\frac{2\sqrt3+3}{2}
\]
\end{solution}

\begin{problem}
甲乙二人在河的南岸 \(O\) 处，隔河在正北方向有一建筑物 \(P\). 甲向正东，乙向正西沿河岸而行，甲每分钟比乙多走 \(a\text{ 米}\). \(10\) 分钟后，甲望建筑物 \(P\) 在北 \(\alpha^\circ\) 西（即北偏西 \(\alpha^\circ\)），乙望建筑物 \(P\) 在北 \(\beta^\circ\) 东（即北偏东 \(\beta^\circ\)），求 \(O\) 与 \(P\) 之间的距离.
\end{problem}

\begin{solution}
设十分钟后甲、乙所在位置分别为 \(A\)、\(B\)，设 \(OB=x\)，则乙向正西走了 \(x\) 米，甲向正东走了 \(x+10a\) 米，所以 \(OA=x+10a\).

\solutionfigure{\bitmapfigure[max width=0.42\linewidth]{img/1964/national/q2_solution_fig1.png}}

设 \(OP=h\). 过 \(A\)、\(B\) 分别作与 \(OP\) 同向的正北方向直线，则由北偏角的定义，\(\alpha\)、\(\beta\) 分别是 \(AP\)、\(BP\) 与这些正北方向直线的夹角. 因此在两个直角三角形中分别有
\(\tan\alpha=\frac{OA}{OP}=\frac{x+10a}{h}\)，\(\tan\beta=\frac{OB}{OP}=\frac{x}{h}\)，即 \(h\tan\alpha=x+10a\)，\(h\tan\beta=x\).
两式相减，得到
\[
h\left( \tan\alpha-\tan\beta \right)=10a
\]
所以
\[
OP=h=\frac{10a}{\tan\alpha-\tan\beta}
\]
\end{solution}

\begin{problem}
解方程 \(x^4+1=0\)；并且证明：平面内表示这个方程的根的四个点是一个正方形的顶点.
\end{problem}

\begin{solution}
将方程左边配方并作平方差分解：
\[
x^4+1=\left( x^2+1 \right)^2-2x^2
=\left( x^2+\sqrt2x+1 \right)\left( x^2-\sqrt2x+1 \right)
\]
因此原方程等价于
\[
\left( x^2+\sqrt2x+1 \right)\left( x^2-\sqrt2x+1 \right)=0
\]
由求根公式，两个二次方程的根分别为
\[
x=\frac{-\sqrt2\pm\sqrt{-2}}{2}=\frac{\sqrt2}{2}\left( -1\pm\i \right)
\]
以及
\[
x=\frac{\sqrt2\pm\sqrt{-2}}{2}=\frac{\sqrt2}{2}\left( 1\pm\i \right)
\]
所以四个根可按逆时针顺序记为 \(x_1=\frac{\sqrt2}{2}\left( 1+\i \right)\)，\(x_2=\frac{\sqrt2}{2}\left( -1+\i \right)\)，\(x_3=\frac{\sqrt2}{2}\left( -1-\i \right)\)，\(x_4=\frac{\sqrt2}{2}\left( 1-\i \right)\).
令 \(c=\frac{\sqrt2}{2}\)，并在复平面内分别以 \(M_1\)，\(M_2\)，\(M_3\)，\(M_4\) 表示这四个根，则
\(M_1=(c,c)\)，\(M_2=(-c,c)\)，\(M_3=(-c,-c)\)，\(M_4=(c,-c)\).
于是
\[
M_1M_2=M_2M_3=M_3M_4=M_4M_1=2c=\sqrt2
\]
且 \(M_1M_2\) 平行于实轴，\(M_2M_3\) 平行于虚轴，所以 \(\angle M_1M_2M_3=90^\circ\). 四边相等且有一个直角，故 \(M_1M_2M_3M_4\) 是正方形.

\solutionfigure{\bitmapfigure[max width=0.38\linewidth]{img/1964/national/q3_solution_fig1.png}}
\end{solution}

\begin{problem}
已知 \(A\)，\(B\)，\(C\) 是三角形的三个内角，求证：\(\cos A=\frac{\sin^2B+\sin^2C-\sin^2A}{2\sin B\sin C}\).
\end{problem}

\begin{solution}
设三角形的三边分别为 \(a\)、\(b\)、\(c\)，其中 \(a\)、\(b\)、\(c\) 分别与角 \(A\)、\(B\)、\(C\) 对应，外接圆半径为 \(R\). 由正弦定理，
\[
\frac{a}{\sin A}=\frac{b}{\sin B}=\frac{c}{\sin C}=2R
\]
即 \(\sin A=\frac{a}{2R}\)，\(\sin B=\frac{b}{2R}\)，\(\sin C=\frac{c}{2R}\).
代入右端，得到
\[
\frac{\sin^2B+\sin^2C-\sin^2A}{2\sin B\sin C}
=\frac{\frac{b^2}{4R^2}+\frac{c^2}{4R^2}-\frac{a^2}{4R^2}}{2\cdot\frac{b}{2R}\cdot\frac{c}{2R}}
=\frac{b^2+c^2-a^2}{2bc}
\]
由余弦定理 \(a^2=b^2+c^2-2bc\cos A\)，所以
\[
\frac{b^2+c^2-a^2}{2bc}=\cos A
\]
又因为 \(B\)、\(C\) 是三角形内角，\(\sin B\sin C>0\)，上述除法合法，故原等式成立.
\end{solution}

\begin{problem}
已知方程 \(x^3+mx^2-3x+n=0\) 的三个根的平方和为 \(6\)，且知这个方程有两相等的正根，求 \(m\)，\(n\) 的值.
\end{problem}

\begin{solution}
设两个相等的正根为 \(\alpha\)、\(\alpha\)，另一个根为 \(\beta\)，其中 \(\alpha>0\). 由三个根的平方和为 \(6\)，得
\[
2\alpha^2+\beta^2=6
\]
由根与系数的关系，得
\[
\alpha^2+2\alpha\beta=-3
\]
将第二式乘以 \(2\)，再与第一式相加，得到
\[
\left( 2\alpha^2+\beta^2 \right)+2\left( \alpha^2+2\alpha\beta \right)=6-6=0
\]
即
\[
\left( 2\alpha+\beta \right)^2=0
\]
因此 \(\beta=-2\alpha\). 代入平方和条件，得 \(2\alpha^2+\left( -2\alpha \right)^2=6\)，即 \(6\alpha^2=6\).
由于 \(\alpha>0\)，所以 \(\alpha=1\)，进而 \(\beta=-2\).

再次由根与系数的关系，\(2\alpha+\beta=-m\)，\(\alpha^2\beta=-n\).
代入 \(\alpha=1\)、\(\beta=-2\)，得
\(m=0\)，\(n=2\).
校验可得方程为 \(x^3-3x+2=\left( x-1 \right)^2\left( x+2 \right)=0\)，确有两个相等的正根 \(1\)，故 \(m=0\)，\(n=2\).
\end{solution}

\begin{problem}
圆台形铁桶的上口半径是 \(15\text{ 厘米}\)，下底半径是 \(10\text{ 厘米}\)，母线长是 \(30\text{ 厘米}\)，将铁桶的侧面沿一条母线剪开铺平，得图中扇面形状的铁片 \(ABCD\). 求 \(A\)，\(B\) 两点间的距离.

\bitmapfigure[max width=0.42\linewidth]{img/1964/national/q6_fig1.png}
\end{problem}

\begin{solution}
设扇形的圆心为 \(O\)，延长 \(AD\) 与 \(BC\) 相交于 \(O\)，设 \(OD=x\text{ 厘米}\)，设 \(\angle COD=\theta\)，其中 \(\theta\) 用弧度表示. 展开后，外侧弧 \(\widehat{AB}\) 的长度等于铁桶上口周长，内侧弧 \(\widehat{CD}\) 的长度等于铁桶下底周长，因此由圆周长公式，\(\widehat{AB}=2\pi\cdot15=30\pi\text{ 厘米}\)，\(\widehat{CD}=2\pi\cdot10=20\pi\text{ 厘米}\).
由弧长公式，得到
\[
\begin{cases}
\left( x+30 \right)\theta=30\pi,\\
x\theta=20\pi.
\end{cases}
\]
两式相减，得 \(30\theta=10\pi\)，所以 \(\theta=\frac{\pi}{3}\). 再由 \(x\theta=20\pi\)，得
\[
x=\frac{20\pi}{\theta}=60
\]
于是 \(OA=OB=x+30=90\text{ 厘米}\)，且 \(\angle AOB=\theta=\frac{\pi}{3}\). 因而
\[
AB=2\cdot90\sin\frac{\theta}{2}=180\sin\frac{\pi}{6}=90\text{ 厘米}
\]
\end{solution}

\begin{problem}
\(A\)，\(B\)，\(C\)，\(D\) 四个点在平面 \(M\) 和平面 \(N\) 之外，\(A\)，\(B\)，\(C\)，\(D\) 在平面 \(M\) 内的射影是 \(A_1\)，\(B_1\)，\(C_1\)，\(D_1\)，在平面 \(N\) 内的射影是 \(A_2\)，\(B_2\)，\(C_2\)，\(D_2\). 已知 \(A_1\)，\(B_1\)，\(C_1\)，\(D_1\) 在一条直线上，\(A_2B_2C_2D_2\) 是一个平行四边形，求证 \(ABCD\) 也是一个平行四边形.
\end{problem}

\begin{solution}
设 \(A_1\)、\(B_1\)、\(C_1\)、\(D_1\) 所在的直线为 \(l\). 过 \(l\) 和直线 \(AA_1\) 作平面 \(P\). 因为 \(l\subset M\)，且 \(AA_1\perp M\)，所以 \(P\perp M\). 显然 \(A\in P\). 对于点 \(B\)，由于 \(B_1\in l\subset P\)，且 \(BB_1\perp M\)，而平面 \(P\perp M\)，所以平面 \(P\) 内过 \(B_1\) 的垂线就是 \(BB_1\)，故 \(BB_1\subset P\)，从而 \(B\in P\). 同理，\(C\)、\(D\) 也在平面 \(P\) 内.

\solutionfigure{\bitmapfigure[max width=0.5\linewidth]{img/1964/national/q7_solution_fig1.png}}

因为 \(AA_2\)、\(DD_2\) 都垂直于平面 \(N\)，所以 \(AA_2\parallel DD_2\). 又因为 \(A_2B_2C_2D_2\) 是平行四边形，所以 \(A_2B_2\parallel D_2C_2\). 设由直线 \(AA_2\)、\(A_2B_2\) 所确定的平面为 \(U\)，由直线 \(DD_2\)、\(D_2C_2\) 所确定的平面为 \(V\). 由于 \(A_2B_2\) 与 \(D_2C_2\) 是不重合的平行直线，且 \(AA_2\parallel DD_2\)，所以 \(U\) 与 \(V\) 是不重合的平行平面.

由于 \(BB_2\parallel AA_2\)，且 \(B_2\in A_2B_2\)，有 \(B\in U\). 若 \(U=P\)，则 \(U\) 含有垂直于 \(N\) 的直线 \(AA_2\)，因而 \(U\perp N\)，且 \(U\cap N=A_2B_2\). 对于 \(C\)，\(D\in P=U\)，它们在 \(N\) 内的正射影 \(C_2\)，\(D_2\) 必在 \(U\cap N=A_2B_2\) 上，这与 \(A_2B_2C_2D_2\) 是平行四边形矛盾，故 \(U\ne P\). 同理 \(V\ne P\). 因此 \(U\cap P=AB\)，\(V\cap P=CD\). 平面 \(P\) 截两个平行平面所得的交线平行，因此 \(AB\parallel CD\).

同样，由 \(AA_2\parallel BB_2\) 和 \(A_2D_2\parallel B_2C_2\)，由 \(AA_2\)、\(A_2D_2\) 所确定的平面与由 \(BB_2\)、\(B_2C_2\) 所确定的平面是两不重合的平行平面. 这两个平面与 \(P\) 的交线分别为 \(AD\)、\(BC\)，所以 \(AD\parallel BC\). 两组对边分别平行，故四边形 \(ABCD\) 是平行四边形.
\end{solution}

\begin{problem}
下图中 \(ABCD\) 是正方形，其每边长为 \(1\)；在正方形内，\(\odot O\) 与 \(\odot O'\) 互相外切，并且 \(\odot O\) 与 \(AB\)，\(AD\) 两边相切，\(\odot O'\) 与 \(CB\)，\(CD\) 两边相切.

\begin{enumerate}
\item 求这两圆半径之和.
\item 当两圆半径各多么长时，两圆面积之和最小？当半径各多么长时，面积之和最大？证明你的结论.
\end{enumerate}

\bitmapfigure[max width=0.32\linewidth]{img/1964/national/q8_fig1.png}
\end{problem}

\begin{solution}
设两圆半径分别为 \(r\)、\(r'\)，令 \(s=r+r'\). 建立直角坐标系，取 \(A=(0,0)\)、\(B=(1,0)\)、\(D=(0,1)\)，则 \(C=(1,1)\). 由于 \(\odot O\) 与 \(AB\)、\(AD\) 相切，\(O=(r,r)\)；由于 \(\odot O'\) 与 \(CB\)、\(CD\) 相切，\(O'=(1-r',1-r')\).

\solutionfigure{\bitmapfigure[max width=0.4\linewidth]{img/1964/national/q8_solution_fig1.png}}

由于两圆都在正方形内且分别与两边相切，\(0<r\)，\(r'\le\frac12\). 若 \(s=1\)，则 \(r=r'=\frac12\)，从而 \(O=O'=\left(\frac12,\frac12\right)\)，这与两圆外切矛盾，故 \(0<s<1\).

\begin{enumerate}
\item 两圆外切，所以 \(OO'=r+r'=s\). 另一方面，
\[
OO'=\sqrt{\left( 1-r'-r \right)^2+\left( 1-r'-r \right)^2}=\sqrt2\left( 1-s \right)
\]
因此 \(s=\sqrt2\left( 1-s \right)\)，解得
\[
s=\frac{\sqrt2}{1+\sqrt2}=2-\sqrt2
\]

\item 每个圆都在边长为 \(1\) 的正方形内并与两边相切，所以 \(0<r\le\frac12\)，\(0<r'\le\frac12\)，并且 \(r+r'=s=2-\sqrt2\).
两圆面积之和为
\[
S=\pi\left( r^2+r'^2 \right)=\frac{\pi}{2}\left[ s^2+\left( r-r' \right)^2 \right]
\]
当且仅当 \(r-r'=0\) 时，\(S\) 取得最小值，所以
\[
r=r'=\frac{s}{2}=\frac{2-\sqrt2}{2}
\]

由 \(0<r\)，\(r'\le\frac12\) 及 \(r+r'=s\)，可知
\[
s-\frac12\le r\le\frac12
\]
在此区间上，\(\left( r-r' \right)^2=\left( 2r-s \right)^2\) 在端点处取得最大值. 因而面积之和最大时，一圆半径为 \(\frac12\)，另一圆半径为
\[
s-\frac12=\frac32-\sqrt2
\]
所以最大值在
\[
\left( r,r' \right)=\left( \frac12,\frac32-\sqrt2 \right)
\]
或两圆互换这两个半径时取得.
\end{enumerate}
\end{solution}

\section{附加题}

\begin{problem}
\begin{enumerate}
\item 如果把第 \(8\) 题中的正方形改成矩形，你能得什么结果？为什么？

\item 如果把第 \(8\) 题中的正方形改成棱长为 \(1\) 的正方体，把圆改成球，你能得到什么结果？为什么？
\end{enumerate}
\end{problem}

\begin{solution}
\begin{enumerate}
\item
设矩形 \(ABCD\) 中 \(AB=a\)、\(BC=b\)，其中 \(a>0\)、\(b>0\). 设两圆半径分别为 \(r\)、\(r'\)，且令 \(s=r+r'\). 取 \(A=(0,0)\)、\(B=(a,0)\)、\(D=(0,b)\)，则两圆心分别为 \(O=(r,r)\) 和 \(O'=(a-r',b-r')\).

\solutionfigure{\bitmapfigure[max width=0.5\linewidth]{img/1964/national/q8_additional_solution_fig1.png}}

两圆外切，所以 \(OO'=s\). 由两圆心坐标，
\[
s^2=OO'^2=\left( a-r-r' \right)^2+\left( b-r-r' \right)^2
=\left( a-s \right)^2+\left( b-s \right)^2
\]
整理得
\[
s^2-2\left( a+b \right)s+a^2+b^2=0
\]
解这个关于 \(s\) 的一元二次方程，得
\[
s=a+b\pm\sqrt{2ab}
\]
令 \(m=\min\{a,b\}\)、\(M=\max\{a,b\}\). 因为每个圆都在矩形内并分别与相邻两边相切，所以 \(0<r\le\frac{m}{2}\)，\(0<r'\le\frac{m}{2}\)，从而 \(0<s\le m\). 因此只能取较小的根：
\[
s=a+b-\sqrt{2ab}
\]
这个构型存在还必须满足
\[
a+b-\sqrt{2ab}\le m
\]
由于 \(a+b=M+m\)，上式等价于
\[
M\le\sqrt{2Mm}\iff M^2\le2Mm\iff M\le2m
\]
若 \(M>2m\)，则上述两圆构型不存在. 反之，当 \(M\le2m\) 时，取 \(r=r'=\frac{s}{2}\le\frac m2\)，两圆均在矩形内且满足外切条件，所以该条件也是充分的. 以下均在 \(M\le2m\) 的条件下讨论. 在此条件下，\(s-\frac m2>0\)，因为
\[
\left( M+\frac m2 \right)^2-2Mm=\left( M-\frac m2 \right)^2>0
\]

两圆面积之和为
\[
S=\pi\left( r^2+r'^2 \right)=\frac{\pi}{2}\left[ s^2+\left( r-r' \right)^2 \right]
\]
由于 \(s\) 是定值，当且仅当 \(r=r'=\frac{s}{2}\) 时面积之和最小，即
\[
r=r'=\frac{a+b-\sqrt{2ab}}{2}
\]

又因为 \(r\)，\(r'\le\frac{m}{2}\) 且 \(r+r'=s\)，可行范围为 \(s-\frac m2\le r\le\frac m2\)，且 \(r'=s-r\). 此时 \(r-r'=2r-s\) 随 \(r\) 增大而增大，所以 \(\lvert r-r'\rvert\) 在区间端点处取得最大值. 因而面积之和最大时，两圆半径为
\[
\left\{r,r'\right\}=\left\{\frac{m}{2},\ a+b-\sqrt{2ab}-\frac{m}{2}\right\}
\]
对应的最大面积之和为
\[
S_{\max}=\pi\left[\left( \frac{m}{2} \right)^2+\left( a+b-\sqrt{2ab}-\frac{m}{2} \right)^2\right]
\]

\item
设两个球的半径分别为 \(r\)、\(r'\)，令 \(s=r+r'\). 建立直角坐标系，使正方体为 \([0,1]^3\)，并使第一个球与过顶点 \((0,0,0)\) 的三个面相切，第二个球与相对顶点 \((1,1,1)\) 的三个面相切. 则两球心分别为 \(O=(r,r,r)\) 和 \(O'=(1-r',1-r',1-r')\).

\solutionfigure{\bitmapfigure[max width=0.34\linewidth]{img/1964/national/q8_additional_solution_fig3.png}}

由于两球都在正方体内且分别与三个面相切，\(0<r\)，\(r'\le\frac12\). 若 \(s=1\)，则 \(r=r'=\frac12\)，从而 \(O=O'=\left(\frac12,\frac12,\frac12\right)\)，这与两球外切矛盾，故 \(0<s<1\).

两球外切，故
\[
s^2=OO'^2=3\left( 1-r-r' \right)^2=3\left( 1-s \right)^2
\]
由于 \(0<s<1\)，取正平方根得到 \(s=\sqrt3\left( 1-s \right)\)，解得
\[
s=\frac{\sqrt3}{1+\sqrt3}=\frac{3-\sqrt3}{2}
\]

由于 \(s=\frac{3-\sqrt3}{2}\) 且 \(\frac12<s<1\)，可行范围为
\(s-\frac12\le r\le\frac12\)，\(r'=s-r\).

两个球的表面积之和为
\[
S=4\pi\left( r^2+r'^2 \right)=2\pi\left[ s^2+\left( r-r' \right)^2 \right]
\]
因为 \(0<r\)，\(r'\le\frac12\) 且 \(r+r'=s\)，当且仅当
\[
r=r'=\frac{s}{2}=\frac{3-\sqrt3}{4}
\]
时表面积之和最小. 当一个半径取允许的最大值 \(\frac12\) 时，另一个半径为
\[
s-\frac12=\frac{2-\sqrt3}{2}
\]
在上述可行区间上，\(\lvert r-r'\rvert=\lvert 2r-s\rvert\) 在端点处取得最大值. 此时表面积之和最大. 也就是说最大值对应于
\[
\left( r,r' \right)=\left( \frac12,\frac{2-\sqrt3}{2} \right)
\]
或两球互换这两个半径.

两个球的体积之和为
\[
V=\frac43\pi\left( r^3+r'^3 \right)
=\frac43\pi s\left( r^2-rr'+r'^2 \right)
=\frac{\pi s}{3}\left[ s^2+3\left( r-r' \right)^2 \right]
\]
其中用到了
\[
4\left( r^2-rr'+r'^2 \right)=s^2+3\left( r-r' \right)^2
\]
这里 \(s\) 为定值，所以体积之和也在 \(r=r'\) 时最小，在一个半径为 \(\frac12\)、另一个半径为 \(\frac{2-\sqrt3}{2}\) 时最大.
\end{enumerate}
\end{solution}
